
\section{Indução e recursão em relações bem-fundadas e set-like}
Nessa seção, estudaremos os princípios de indução e recursão em relações bem-fundadas e set-like.
Ao final, apresentaremos a noção de $R$-rank de um conjunto, em que $R$ é uma relação definida set-like e bem-fundada.

\subsection{Indução transfinita}

Iniciaremos discutindo um princípio de minimalidade e um princípio de indução para relações definidas set-like e bem-fundadas.

Pensamos em $a R b$ como ``$a$ é um $R$-predecessor de $b$'', ou ``$a$ $R$-precede $b$''.

A proposição abaixo é uma generalização do princípio da boa ordenação do conjunto dos números naturais e da classe dos números ordinais.
Ela diz que, dada uma fórmula $\phi(x, t_1, \dots, t_n)$, fixados parâmetros $t_1, \dots, t_n$, se existe um $x$ que satisfaz a fórmula $\phi(x, t_1, \dots, t_n)$, então existe um tal $x$ que é $R$-minimal nesse sentido: ou seja, nenhum $y$ que o $R$-precede satisfaz $\phi(y, t_1, \dots, t_n)$.

\begin{proposition}[$\ZFmP$](Princípio da minimalidade)\label{prop:principioDaMinimalidade}
    Seja $R$ uma relação definida set-like e bem-fundada.
    
    Seja $\phi(x, t_1, \dots, t_n)$ uma fórmula.
    Então:
    \[
        \forall t_1\dots \forall t_n\bigg(\exists x \, \phi(x, t_1, \dots, t_n) \rightarrow \exists x\bigg(\phi(x, t_1, \dots, t_n)\wedge \forall y\big(y R x \rightarrow \neg \phi(y, t_1, \dots, t_n)\big)\bigg)\bigg)
    \]
\end{proposition}

\begin{proof}
    Fixe $t_1, \dots, t_n$ e suponha que $\exists x\, \phi(x, t_1, \dots, t_n)$.
    Fixe $x$ dessa forma.

    Se $\forall y\big(y R x \rightarrow \neg \phi(y, t_1, \dots, t_n)\big)$, segue a tese.
    Caso contrário, considere
    \[
        X=\{z \in \pred_{R^*}(x): \phi(z, t_1, \dots, t_n)\}.
    \]
    Temos que $X$ é um conjunto e $X\neq \emptyset$, pois $R^*$ estende $R$.
    Logo, existe um elemento $R$-minimal de $X$, digamos, $x'$.
    Temos que $\phi(x', t_1, \dots, t_n)$.
    Além disso, se $y R x'$, então, como $x' R^* x$ e $R^*$ é transitiva e estende $R$, segue que $y\in \pred_{R^*}(x)$.
    Pela minimalidade de $x'$ em $X$, tal $y$ não pertence a $X$; logo, $\neg \phi(y, t_1, \dots, t_n)$.
\end{proof}

Como corolário, provamos o princípio da indução sobre uma relação definida set-like e bem-fundada $R$.

Ele diz que, dada uma fórmula $\phi(x, t_1, \dots, t_n)$ e fixados os parâmetros $t_1, \dots, t_n$, se, para todo $x$, a validade de $\phi(z, t_1, \dots, t_n)$ em todos os $R$-predecessores $z$ de $x$ implica a validade de $\phi(x, t_1, \dots, t_n)$, então $\phi(x, t_1, \dots, t_n)$ vale para todo $x$.

\begin{corollary}[$\ZFmP$](Princípio da indução)\label{cor:principioDaInducao}
    Seja $R$ uma relação definida set-like e bem-fundada.
    
    Seja $\phi(x, t_1, \dots, t_n)$ uma fórmula.
    Então:
    \[
        \forall t_1\dots \forall t_n\bigg(\forall x \, \left(\forall z \in \pred_{R}(x) \, \phi(z, t_1, \dots, t_n) \rightarrow \, \phi(x, t_1, \dots, t_n)\right)\rightarrow \forall x\, \phi(x, t_1, \dots, t_n)\bigg)
    \]
\end{corollary}

\begin{proof}
    Fixe $t_1, \dots, t_n$ e suponha que existe $x$ tal que $\neg \phi(x, t_1, \dots, t_n)$.

    Pela Proposição~\ref{prop:principioDaMinimalidade}, existe um elemento $R$-minimal $x'$ tal que $\neg \phi(x', t_1, \dots, t_n)$ e $\phi(z, t_1, \dots, t_n)$ para todo $z\in \pred_{R}(x')$.

    Por hipótese, temos que $\forall x \, \left(\forall z \in \pred_{R}(x) \, \phi(z, t_1, \dots, t_n) \rightarrow \, \phi(x, t_1, \dots, t_n)\right)$.
    Aplicando isso para $x'$, temos que $\phi(x', t_1, \dots, t_n)$, o que é uma contradição.
\end{proof}

Ao utilizarmos indução em uma relação definida set-like e bem-fundada $R$, por vezes, escreveremos que faremos uso de uma \emph{$R$-indução}.

\subsection{Recursão transfinita}
Assim como a indução em uma relação set-like e bem-fundada $R$ permite provar propriedades assumindo que elas já valem nos $R$-predecessores, a recursão permite definir objetos assumindo que seus valores já foram definidos nesses predecessores.

Formalmente, precisamos de uma operação $G(p, s, a)$ que recebe um parâmetro $p$, um ponto $a$ e um ``histórico'' $s$ de valores já definidos para os $R$-predecessores de $a$, e retorna o valor a ser definido em $a$.

O Teorema da Recursão, na forma apresentada abaixo, diz que, dada uma tal $G$, é possível definir sua versão iterada $F(p, a)$, que, para cada $a$, retorna o valor definido para $a$ a partir de $G$, do parâmetro $p$ e do histórico $F_p|\pred_{R}(a)$ de valores definidos para os $R$-predecessores de $a$.

\begin{theorem}[$\ZFmP$](Esquema da recursão transfinita)\label{thm:esquemaDaRecursaoTransfinita}\index{recursão transfinita}
    Seja $R$ uma relação definida set-like e bem-fundada.
    Suponha que um símbolo funcional ternário $G(p, s, a)$ seja introduzido por definição.

    Então é possível introduzir um símbolo funcional binário $F(p, a)$ de modo que
    \[
        \forall a \forall p \Big( F_{p}(a)=G(p, F_p|\pred_{R}(a), a)\Big).
    \]
\end{theorem}
Para provar esse teorema, façamos as seguintes definições auxiliares.

\begin{definition}[$\ZFmP$]
    Seja $R$ uma relação definida set-like.
    Seja $G(p, s, a)$ um símbolo funcional ternário introduzido por definição.

    Dados conjuntos $t, a, p$, escrevemos $\Ap_G(t, p, a)$ e dizemos que \emph{$t$ é uma computação parcial em $a$ com parâmetro $p$ baseada em $G$} se, e somente se,
    \begin{itemize}
        \item $t$ é uma função,
        \item $\dom(t)$ é $R$-fechado,
        \item $a \in \dom(t)$,
        \item $t(b)=G(p, t|\pred_{R}(b), b)$ para todo $b\in \dom(t)$.
    \end{itemize}
\end{definition}

\begin{lemma}[$\ZFmP$]
    Seja $R$ uma relação definida set-like e bem-fundada e $G$ um símbolo funcional ternário introduzido por definição.

    Se $t_1$, $t_2$, $a_1$, $a_2$, $p$ são tais que $\Ap_G(t_1, p, a_1)$ e $\Ap_G(t_2, p, a_2)$, então $t_1$ e $t_2$ são compatíveis.
\end{lemma}
\begin{proof}
    Suponha por absurdo que não.
    Então o conjunto
    \[
        X=\{x\in \dom(t_1)\cap \dom(t_2): t_1(x)\neq t_2(x)\}
    \]
    é não vazio.
    Seja $x$ um elemento $R$-minimal de $X$.

    Dessa forma, para todo $y R x$, temos que $y\in \dom(t_1)\cap \dom(t_2)$, pois $\dom(t_1)$ e $\dom(t_2)$ são $R$-fechados, e, portanto, $t_1(y)=t_2(y)$ pela minimalidade de $x$.
    Logo, $t_1|\pred_{R}(x)=t_2|\pred_{R}(x)$.
    Daí, segue que:
    \[
        t_1(x)=G(p, t_1|\pred_{R}(x), x)=G(p, t_2|\pred_{R}(x), x)=t_2(x),
    \]
    o que é uma contradição.
\end{proof}

O lema anterior garante a unicidade das computações parciais, pois ele diz que duas computações parciais construídas pela mesma regra recursiva nunca entram em conflito nos pontos em que ambas estão definidas.
Por isso, mais adiante, quando reunirmos várias computações parciais, a união obtida continuará sendo uma função.

\begin{corollary}[$\ZFmP$]
    Seja $R$ uma relação definida set-like e bem-fundada e $G$ um símbolo funcional ternário introduzido por definição.

    Sejam $a$ e $p$ conjuntos.
    Se existe $t$ tal que $\Ap_G(t, p, a)$, então existe um único $t'$ tal que $\Ap_G(t', p, a)$ e $\dom(t')=\pred_{R^*}(a)\cup\{a\}$.
\end{corollary}
\begin{proof}
    Para a existência, fixado $t$, tome $t'=t|(\pred_{R^*}(a)\cup\{a\})$.
    Como $\dom(t)$ é $R$-fechado e $a\in \dom(t)$, temos $\pred_R(a)\subseteq \dom(t)$.
    Pelo Lema~\ref{lem:predRStarFechado} aplicado à $\dom(t)$, temos $\pred_{R^*}(a)\subseteq \dom(t)$.
    Logo, $\pred_{R^*}(a)\cup\{a\}\subseteq \dom(t)$.

    Além disso, $\dom(t')=\pred_{R^*}(a)\cup\{a\}$ é $R$-fechado.
    De fato, se $b=a$, então $\pred_R(b)=\pred_R(a)\subseteq \pred_{R^*}(a)\subseteq \dom(t')$, e, se $b\in \pred_{R^*}(a)$, então $\pred_R(b)\subseteq \pred_{R^*}(a)$, pois $\pred_{R^*}(a)$ é $R$-fechado.

    Portanto, se $b\in \dom(t')$, então $\pred_R(b)\subseteq \dom(t')$, de modo que $t'|\pred_R(b)=t|\pred_R(b)$.
    Assim, $t'$ satisfaz as mesmas equações recursivas que $t$ em todos os pontos de seu domínio.

    Para a unicidade, se $t'$ e $t''$ são tais que $\Ap_G(t', p, a)$ e $\Ap_G(t'', p, a)$ e $\dom(t')=\dom(t'')=\pred_{R^*}(a)\cup\{a\}$, então $t'$ e $t''$ são compatíveis, e, portanto, $t'=t''$.
\end{proof}
\begin{lemma}[$\ZFmP$]
    Seja $R$ uma relação definida set-like e bem-fundada e $G$ um símbolo funcional ternário introduzido por definição.

    Para todo $a$ e $p$, existe um $t$ tal que $\Ap_G(t, p, a)$.
\end{lemma}

\begin{proof}
    Fixe $p$.
    Provaremos por $R$-indução que, para todo $a$, existe um $t$ tal que $\Ap_G(t, p, a)$.

    Dado $a$, basta ver que se a tese vale para todo $b\in \pred_{R}(a)$, então vale para $a$.

    Suponha que, para todo $c\in \pred_{R}(a)$, existe um $t$ tal que $\Ap_G(t, p, c)$.
    Pelo corolário anterior, para cada $c\in \pred_{R}(a)$, seja $t_c$ o único $t$ tal que $\Ap_G(t, p, c)$ e $\dom(t)=\pred_{R^*}(c)\cup\{c\}$.
    Pela Substituição, podemos considerar a família $\{t_c: c\in \pred_{R}(a)\}$.

    Pelo lema de compatibilidade, os $t_c$ são dois a dois compatíveis.
    Portanto,
    \[
        t'=\bigcup_{c\in \pred_{R}(a)} t_c
    \]
    é uma função.
    Além disso,
    \[
        \dom(t')=\bigcup_{c\in \pred_{R}(a)} \dom(t_c)
        =\bigcup_{c\in \pred_{R}(a)}(\pred_{R^*}(c)\cup\{c\})
        =\pred_{R^*}(a).
    \]

    Como $R$ é bem-fundada, temos $a\notin \pred_{R^*}(a)$; caso contrário, os elementos de um $R$-caminho de $a$ até $a$ formariam um conjunto não vazio sem elemento $R$-minimal.
    Seja
    \[
        t=t'\cup\{(a, G(p, t'|\pred_{R}(a), a))\}.
    \]
    Então $t$ é função, $\dom(t)=\pred_{R^*}(a)\cup\{a\}$, esse domínio é $R$-fechado e $a\in \dom(t)$.

    Dado $b\in \dom(t)$, se $b=a$, então $\pred_{R}(a)\subseteq \dom(t')$, de modo que $t|\pred_{R}(a)=t'|\pred_{R}(a)$ e
    \[
        t(a)=G(p, t'|\pred_{R}(a), a)=G(p, t|\pred_{R}(a), a).
    \]

    Já se $b\neq a$, existe um $c\in \pred_{R}(a)$ tal que $b\in \pred_{R^*}(c)\cup\{c\}$.
    Como $\dom(t_c)$ é $R$-fechado, $\pred_{R}(b)\subseteq \dom(t_c)$.
    Além disso, $t$ e $t_c$ coincidem em $\dom(t_c)$.
    Portanto,
    \[
        t(b)=t_c(b)=G(p, t_c|\pred_{R}(b), b)=G(p, t|\pred_{R}(b), b).
    \]
\end{proof}

Abaixo, re-enunciamos o Teorema~\ref{thm:esquemaDaRecursaoTransfinita} de maneira mais precisa.

\begin{theorem}[$\ZFmP$](Esquema da recursão transfinita)
    Seja $R$ uma relação definida set-like e bem-fundada.
    Suponha que um símbolo funcional ternário $G(p, s, a)$ seja introduzido por definição.

    Seja $F^G$ o símbolo funcional binário introduzido por definição tal que:
    \[
        \forall a \forall p \forall y\bigg(F^G_{p}(a)=y \leftrightarrow \exists t\big(\Ap_G(t, p, a)\wedge t(a)=y\big)\bigg).
    \]
    Então $\forall a \forall p \Big( F^G_{p}(a)=G(p, F^G_p|\pred_{R}(a), a)\Big)$.
\end{theorem}
\begin{proof}
    Primeiro, notemos que $F^G$ é bem definida: dado $a$ e $p$, existe um $t$ tal que $\Ap_G(t, p, a)$, e, se $t'$ é outro tal que $\Ap_G(t', p, a)$, então $t'$ e $t$ são compatíveis, e, portanto, $t'(a)=t(a)$.

    Segundo, dado $a$ e $p$, seja $t$ tal que $\Ap_G(t, p, a)$.
    Então, para todo $b\in \pred_{R}(a)$, temos $b\in \dom(t)$ e, portanto, $\Ap_G(t, p, b)$.
    Pela definição de $F^G$, segue que $t|\pred_{R}(a)=F^G_p|\pred_{R}(a)$.
    Assim, $F^G_p(a)=t(a)=G(p, t|\pred_{R}(a), a)=G(p, F^G_p|\pred_{R}(a), a)$.
\end{proof}

\subsection{Ranks de relações bem-fundadas e set-like}

De posse da recursão transfinita, podemos definir o rank de um conjunto com respeito a uma relação definida set-like e bem-fundada.

Intuitivamente, $\rank_R(x)$ mede a ``altura'' de $x$ na relação $R$: elementos sem $R$-predecessores terão rank $0$, e o rank de $x$ será sempre estritamente maior do que o rank de cada $R$-predecessor de $x$.
\begin{definition}[$\ZFmP$]
    Seja $R$ uma relação definida set-like e bem-fundada.
    O \emph{rank} de um conjunto $x$ com respeito a $R$, denotado por $\rank_R(x)$, é definido recursivamente pela relação $R$ a partir da seguinte fórmula

    \[
        \rank_R(x)=\bigcup\{S(\rank_R(y)): y R x\},
    \]
    em que $S(u)=u\cup\{u\}$ é a operação sucessora.
\end{definition}

A definição acima pode ser formalizada a partir do teorema da recursão transfinita, tomando
\[
    G(p, s, a)=\bigcup\{S(s(b)): b\in \dom s \land b R a\}.
\]
O Teorema da Recursão nos define um símbolo funcional binário $F(p, a)$ tal que, para todos $p, a$:
\[
    F_p(a)=G(p, F_p|\pred_{R}(a), a)=\bigcup\{S(F_p(b)): b R a\}.
\]

Tome $p$ qualquer, digamos, $p=\emptyset$, e seja $\rank_R(a)=F_p(a)$.
Então $\rank_R(a)=\bigcup\{S(\rank_R(b)): b R a\}$, como queríamos.

Tal função é, em algum sentido, a única função que satisfaz essa fórmula (ver Exercício~\ref{ex:rankBemDefinido}).

\begin{lemma}[$\ZFmP$]
    Seja $R$ uma relação definida set-like e bem-fundada.
    \begin{enumerate}[label=(\roman*)]
        \item Para todo $x$, $\rank_R(x)$ é um ordinal.
        \item Para todo $x$, $\rank_R(x)=\sup\{\rank_R(y)+1: y R x\}$.
        \item Para todos $x$ e $y$, se $y R x$, então $\rank_R(y)<\rank_R(x)$.
        \item Para todos $x$ e $y$, se $y R^* x$, então $\rank_R(y)<\rank_R(x)$.
    \end{enumerate}
\end{lemma}
\begin{proof}
    Para (i), por $R$-indução, para todo $x$, $\rank_R(x)$ é um ordinal.
    Dado $x$, suponha que, para todo $y R x$, $\rank_R(y)$ é um ordinal.
    Então $\rank_R(x)=\bigcup\{S(\rank_R(y)): y R x\}$.
    Como $\rank_R(y)$ é um ordinal para todo $y R x$, temos $S(\rank_R(y))=\rank_R(y)+1$ para todo $y R x$.
    Logo, $\rank_R(x)$ é um ordinal.

    Para (ii), de (i), temos que $\rank_R(x)=\bigcup\{\rank_R(y)+1: y R x\}=\sup\{\rank_R(y)+1: y R x\}$, pois o supremo de um conjunto de ordinais é sua união.

    Para (iii), se $y R x$, temos que $\rank_R(y)+1\leq \rank_R(x)$ por (ii), e, portanto, $\rank_R(y)<\rank_R(x)$.

    Para (iv), suponha que $y R^* x$.
    Então existe um número natural $n\geq 1$ e um $R$-caminho com $n$ passos de $y$ até $x$, digamos, $(x_0,\dots,x_n)$, em que $x_0=y$ e $x_n=x$.
    Por (iii), $\rank_R(x_0)<\rank_R(x_1)<\cdots <\rank_R(x_n)$, e, portanto, $\rank_R(y)=\rank_R(x_0)<\rank_R(x_n)=\rank_R(x)$.
\end{proof}

Adotaremos a seguinte notação, em analogia à Definição~\ref{def:predRA}, para nos referirmos ao rank de um elemento $x$ com respeito a uma relação definida $R$ restrita a uma classe $A$.

\begin{definition}
    Seja $R$ uma relação definida e seja $A$ uma classe tal que $R_A$ é set-like e bem-fundada.
    Para $x\in A$, escrevemos $\rank_{R,A}(x)=\rank_{R_A}(x)$.
\end{definition}

A operação $\rank_R$ fornece uma medida ordinal que cresce estritamente com a relação $R$.
Assim, para provar alguns resultados sobre $R$, por vezes, lança-se mão da tomada de elementos de rank mínimo.
Esse artifício será utilizado na próxima demonstração.

\begin{proposition}[$\ZFmP$]\label{prop:fechoTransitivoBemFundadoEmA}
    Se $A$ é uma classe $R$-fechada e $R$ é uma relação definida set-like e bem-fundada em $A$, então $R^*$ é set-like e é bem-fundada em $A$.
\end{proposition}
\begin{proof}
    Já vimos que $R^*$ é set-like.

    Para ver que $R^*$ é bem-fundada em $A$, seja $X\subseteq A$ um conjunto não vazio.
    Seja $O=\{\rank_{R, A}(x): x\in X\}$.
    Como $O$ é um conjunto de ordinais, existe um ordinal $\alpha$ tal que $\alpha=\min O$.
    Seja $x\in X$ tal que $\rank_{R, A}(x)=\alpha$.
    Se $y\in X$ e $y R^* x$, então, como $A$ é $R$-fechada e $x\in A$, todo $R$-caminho de $y$ até $x$ está contido em $A$.
    Logo, $y R_A^* x$.
    Portanto, $\rank_{R, A}(y)<\rank_{R, A}(x)=\alpha=\min O$, o que contradiz o fato de $\rank_{R,A}(y)\in O$.
    Assim, não existe $y\in X$ tal que $y R^* x$, e, portanto, $x$ é um elemento $R^*$-minimal de $X$.
\end{proof}

Tomando $A=\V$ na proposição anterior, segue que:

\begin{corollary}[$\ZFmP$]
    Se $R$ é uma relação definida set-like e bem-fundada, então $R^*$ é set-like e bem-fundada.
\end{corollary}



\begin{proposition}[$\ZFmP$]\label{prop:relacoesComparaveis2}
    Sejam $R$ e $S$ relações definidas ou relações-conjunto tais que $R\subseteq S$.
    Se $S$ é set-like e bem-fundada, então para todo $x$, $\rank_R(x)\leq \rank_S(x)$.
\end{proposition}
\begin{proof}
    Pela Proposição~\ref{prop:relacoesComparaveis}, $R$ é set-like e bem-fundada.
    Por $R$-indução, veremos que para todo $x$, $\rank_R(x)\leq \rank_S(x)$.
    Dado $x$, suponha que, para todo $y R x$, $\rank_R(y)\leq \rank_S(y)$.
    Então $\rank_R(x)=\sup\{\rank_R(y)+1: y R x\}\leq \sup\{\rank_S(y)+1: y R x\}\leq \sup\{\rank_S(y)+1: y S x\}=\rank_S(x)$.
\end{proof}

\begin{proposition}[$\ZFmP$]\label{prop:rankRestricao}
    Seja $R$ uma relação definida set-like e bem-fundada e $A$ uma classe.
    Se $x\in A$ é tal que $\pred_{R^*}(x)\subseteq A$, então $\rank_{R,A}(x)=\rank_R(x)$.
\end{proposition}
\begin{proof}
    Como $R_A\subseteq R$, a Proposição~\ref{prop:relacoesComparaveis} garante que $R_A$ é set-like e bem-fundada.

    Por $R$-indução, seja dado $x$ tal que a tese vale para todo $y R x$.
    Suponha que $x \in A$ e que $\pred_{R^*}(x)\subseteq A$.

    Dado $y R x$, temos que $y\in A$ e que $\pred_{R^*}(y)\subseteq \pred_{R^*}(x)\subseteq A$.
    Assim, $\rank_{R,A}(y)=\rank_R(y)$.

    Temos, ainda, que para todo $y$, $y R_A x$ se, e somente se, $y R x$, pois se $y R x$, então $y\in \pred_{R^*}(x)\subseteq A$, e, portanto, $y R_A x$.

    Logo, $\rank_{R,A}(x)=\sup\{\rank_{R,A}(y)+1: y R_A x\}=\sup\{\rank_R(y)+1: y R x\}=\rank_R(x)$.
\end{proof}
\section*{Exercícios}

\begin{exercise}[$\ZFmP$]\label{ex:rankBemDefinido}
    Seja $R$ uma relação definida set-like e bem-fundada.
    Suponha que $F$ é um símbolo funcional unário introduzido por definição tal que, para todo $x$, $F(x)=\bigcup\{S(F(y)): y R x\}$.
    Mostre que, para todo $x$,
    \[
        F(x)=\rank_R(x).
    \]
\end{exercise}

\begin{exercise}[$\ZFmP$]\label{exr:relacaoBemFundadaEmVranks}
    Seja $R$ uma relação definida.
    Mostre que se $R$ é set-like e bem-fundada, então $\rank_R(x)=\rank_{R, \V}(x)$ para todo $x$.
\end{exercise}

\begin{exercise}[$\ZFmP$]\label{ex:semCiclosBemFundada}
    Seja $R$ uma relação definida set-like e bem-fundada.

    Mostre que não existem $x$ e $n\geq 1$ tais que existe um $R$-caminho com $n$ passos de $x$ até $x$.

    Conclua que, para todo $x$, não vale $x R^* x$.
\end{exercise}
